\newcommand{\nomeEstado}{BAHIA}
\newcommand{\siglaEstado}{BA}
\newcommand{\nomeConcessionaria}{NEOENERGIA BAHIA}
\newcommand{\cnpjConcessionaria}{15.139.629/0001-94}
\newcommand{\termoCIP}{Sexto}
\newcommand{\presidenteConcessionaria}{Sr. Thiago Guth.}

\newcommand{\empresaResponsavel}{ABEL}
\newcommand{\nomeMunicipio}{Nova Ibiá}
\newcommand{\nomeMunicipioMaisculo}{\MakeUppercase{\nomeMunicipio}}
\newcommand{\IPestimada}{3469025}
\newcommand{\cnpjMunicipio}{13.882.616/0001-04}
\newcommand{\telefone}{(73) 3653-1177}
\newcommand{\email}{prefeitura@novaibia.ba.gov.br}
\newcommand{\sedePrefeitura}{Praça São José, nº 88, Centro, Nova Ibiá – BA, CEP 45452-000}
\newcommand{\nomePrefeito}{Marcio Tarantine Souza Sampaio}
\newcommand{\generoPrefeitura}{pelo Prefeito}
\newcommand{\numeroContrato}{076/2025}
\newcommand{\quantidadeAditivos}{1}
\newcommand{\legitimidade}{1}
\newcommand{\publicacao}{Diário Oficial Eletrônico}
\newcommand{\dataPublicacao}{14 de Março de 2025}
\newcommand{\tituloClausula}{CLÁUSULA PRIMEIRA - DO OBJETO}
\newcommand{\clausula}{[...]

Deverá a empresa efetuar levantamento dos créditos a que faz jus o Município, referentes aos pagamentos indevidos à concessionária de energia elétrica, conforme os prazos estabelecidos na Resolução Normativa da ANEEL nº 1.000, de 7 de dezembro de 2021, especialmente o art. 323, § 2º, inciso I, e suas alterações posteriores, podendo, para tanto, representar administrativamente o Município de Nova Ibiá/BA perante a Agência Nacional de Energia Elétrica — ANEEL, em todos os processos, reclamações, recursos e demais atos decorrentes de suas competências.}
